\section{Recomendaciones}

A partir de los resultados obtenidos, se plantean las siguientes recomendaciones con el propósito de fortalecer la investigación y orientar futuros desarrollos en el área:

\begin{itemize}
    \item \textbf{Ampliar las pruebas de campo:} se recomienda evaluar el sistema en entornos más amplios y con obstáculos, con el fin de estudiar la degradación de la señal UWB y el desempeño del filtro de Kalman en condiciones no controladas.
    
    \item \textbf{Optimizar la calibración de los sensores:} pequeñas variaciones en la altura o el ángulo de los sensores pueden introducir errores sistemáticos en las mediciones. Una rutina de calibración automática previa a cada sesión podría mejorar la precisión global del sistema.
    
    \item \textbf{Explorar métodos de fusión sensorial avanzados:} se sugiere analizar la implementación de filtros de partículas o filtros de Kalman extendidos (EKF), que podrían ofrecer un mejor desempeño en presencia de modelos no lineales y distribuciones no gaussianas.
    
    \item \textbf{Completar y publicar la documentación técnica:} se recomienda consolidar la documentación del plugin, incluyendo diagramas de arquitectura, flujos de datos y guías de integración, de manera que futuras generaciones puedan continuar el desarrollo del megaproyecto sin dificultad.
    
    \item \textbf{Investigar la integración con visión por computadora:} como línea de trabajo futura, podría incorporarse información visual mediante técnicas de SLAM o reconocimiento de marcadores, combinando posicionamiento UWB y percepción visual para mejorar la precisión en zonas con interferencias.
\end{itemize}